\subsection{Bauteilgleichungen}
\setlength{\tabcolsep}{6pt}
\begin{tabular*}{\columnwidth}{@{\extracolsep\fill}lll@{}} \trule
    Widerst. $R = \rho \frac{l}{A}$ & Kondens. $C=\frac{Q}{U} = \varepsilon \frac{A}{d}$ & Spule $L=\mu A \frac{N^2}{l}$\\
    & $W_{el} = \frac{1}{2} C U^2$
\end{tabular*}
\setlength{\tabcolsep}{4pt}



\subsection{Formeln der Elektrostatik}
Coulombsches Gesetz: $\vec F_{q \leftarrow q_i} = \frac{q}{4 \pi \varepsilon} \sum \limits_{i = 1}^N\frac{q_i (\vec r - \vec r_i )}{\abs{\vec r - \vec r_i}^3}$ \\
Elektrische Feldstärke: $\vec E = \frac{\vec F}{q} $ \quad $\rot E = 0$ \\
\textbf{Elektrostatische Felder sind konservativ} $\Leftrightarrow U = \Phi(P_1) - \Phi(P_2) = \int \limits_{P_1}^{P_2} \vec E d \vec r$ ist wegunabhängig \\
Potential: $\Phi (\vec r) = \frac{1}{4 \pi \varepsilon} \sum \limits_{i =1}^{N} \frac{q_i}{\abs{\vec r - \vec r_i}}$ \\
Poissongleichung: $\div(\varepsilon \nabla(\Phi)) = - \rho$ mit $\vec E = -\nabla \Phi$ \\
Oberflächenladungsdichte: $\sigma = \vec D  \vec N$ \\
Energie: $W_{12} = \int_C \vec F d \vec r = q  U_{12}$ \\
Energiedichte: $w_{el} = \frac{1}{2} \vec E \vec D$ \\